\subsection{Base Identity}
Class \texttt{0x00000000}, leaf \texttt{0x0000} reports (:cpuid-field-target:base.cpuid.BASE.IDENTITY.HEADER.MAX_LEAF:)\texttt{MAX\_LEAF}\texttt{=0} and (:cpuid-field-target:base.cpuid.BASE.IDENTITY.HEADER.MAX_INDEX:)\texttt{MAX\_INDEX}\texttt{=18} in its
index-zero (:term:base.terms.header:). An implementation recognizing only the classes defined by this specification reports
(:cpuid-field-target:base.cpuid.BASE.IDENTITY.HEADER.MAX_CLASS:)\texttt{MAX\_CLASS}\texttt{=2}; an implementation that recognizes a higher implementation-defined class reports the highest
such selector. The base-identity (:term:base.terms.payload:) is:

\begin{BedrockListedFormatDiagram}{CPUID Base Identity Result Formats}
\BedrockFormatRowRange{index 1}{63}{0}{%
\BedrockFormatField{VENDOR\_ID}{16}
\BedrockFormatField{ARCHITECTURE\_ID}{16}
\BedrockFormatField{IMPLEMENTATION\_ID}{32}
}
\BedrockFormatRowRange{index 2}{63}{0}{%
\BedrockFormatField{zero}{32}
\BedrockFormatField{IMPL\_REV}{16}
\BedrockFormatField{ARCH\_REV}{16}
}
\end{BedrockListedFormatDiagram}

\BedrockTableCaption{Base Identity Indexes}
\begin{BedrockTabular}{@{}>{\raggedright\arraybackslash}p{0.65in}>{\raggedright\arraybackslash}p{4.95in}@{}}
\toprule
\rowcolor{BedrockHeaderFill}
\textbf{Index} & \textbf{Contents}\\
\midrule
1 & bits 31..0 (:cpuid-field-target:base.cpuid.BASE.IDENTITY.IDENTITY.IMPLEMENTATION_ID:)\texttt{IMPLEMENTATION\_ID}; bits 47..32 (:cpuid-field-target:base.cpuid.BASE.IDENTITY.IDENTITY.ARCHITECTURE_ID:)\texttt{ARCHITECTURE\_ID}; bits 63..48 (:cpuid-field-target:base.cpuid.BASE.IDENTITY.IDENTITY.VENDOR_ID:)\texttt{VENDOR\_ID}\\
2 & bits 15..0 (:cpuid-field-target:base.cpuid.BASE.IDENTITY.REVISIONS.ARCHITECTURE_REVISION:)\texttt{ARCHITECTURE\_REVISION}; bits 31..16 (:cpuid-field-target:base.cpuid.BASE.IDENTITY.REVISIONS.IMPLEMENTATION_REVISION:)\texttt{IMPLEMENTATION\_REVISION}; bits 63..32 zero\\
3--10 & 64-byte vendor name\\
11--18 & 64-byte processor name\\
\bottomrule
\end{BedrockTabular}

(:ref:base.cpuid.BASE.IDENTITY.IDENTITY.VENDOR_ID:) selects the vendor namespace in which (:ref:base.cpuid.BASE.IDENTITY.IDENTITY.IMPLEMENTATION_ID:) is interpreted.
(:ref:base.cpuid.BASE.IDENTITY.IDENTITY.IMPLEMENTATION_ID:) is assigned within that (:ref:base.cpuid.BASE.IDENTITY.IDENTITY.VENDOR_ID:) namespace.
(:ref:base.cpuid.BASE.IDENTITY.IDENTITY.ARCHITECTURE_ID:) identifies the base ISA definition implemented by the processor, while the
((:ref:base.cpuid.BASE.IDENTITY.IDENTITY.VENDOR_ID:), (:ref:base.cpuid.BASE.IDENTITY.IDENTITY.IMPLEMENTATION_ID:)) pair identifies the implementation.

(:ref:base.cpuid.BASE.IDENTITY.REVISIONS.ARCHITECTURE_REVISION:) is the unsigned 16-bit revision of the base ISA definition selected by
(:ref:base.cpuid.BASE.IDENTITY.IDENTITY.ARCHITECTURE_ID:); it does not imply the behavior of any other revision. Optional instruction groups,
(:term:base.terms.processor_state:), and feature-specific parameters are reported separately by their CPUID leaves and feature bits.
(:ref:base.cpuid.BASE.IDENTITY.REVISIONS.IMPLEMENTATION_REVISION:) is an unsigned revision assigned
by the vendor and is compared only when both (:ref:base.cpuid.BASE.IDENTITY.IDENTITY.VENDOR_ID:) and (:ref:base.cpuid.BASE.IDENTITY.IDENTITY.IMPLEMENTATION_ID:) match.

The name (:term:base.terms.field|plural:) are UTF-8 byte strings in increasing index order. Within each 64-bit result, the first byte occupies
bits 7..0 and subsequent bytes occupy successively higher bits. A name shorter than 64 bytes is terminated by a zero
byte and all remaining bytes are zero. A 64-byte name has no terminator. Producers must not split a multi-byte UTF-8
code point at the (:term:base.terms.field:) boundary. Consumers stop at that boundary and replace any invalid sequence.
